\begin{definition}[Left-canonical periodic tensor]
    \label{def:periodic_tensor}
    \lean{MPSTensor.IsPeriodic}
    \leanok
    A left-canonical irreducible MPS tensor $A$ of bond dimension $D$ is
    \emph{$m$-periodic} (for some $m \ge 1$) if the peripheral spectrum of
    its transfer map $\E_A$ is exactly the $m$-th roots of unity
    $\{e^{2\pi i k/m} : 0 \le k < m\}$. This is the normalized form-II
    convention used later in the source.
    A $1$-periodic tensor is exactly an irreducible tensor with
    primitive transfer map (Definition~\ref{def:primitive_channel}).
\end{definition}

\begin{definition}[Irreducible form]
    \label{def:irreducible_form}
    \lean{MPSTensor.IsIrreducibleForm}
    \leanok
    A tensor $A$ is in \emph{irreducible form} if it is block-diagonal
    \[
        A^i = \bigoplus_{j \in J} \mu_j\, A_j^i,
    \]
    where each block $A_j$ is an $m_j$-periodic tensor. Following
    \cite[Section~II]{DeLasCuevas2017Irreducible}, the phases are fixed by the
    convention $\mu_j>0$, so the weights are positive real scalars.
    % Source check: arXiv:1708.00029, lines 252--261 define irreducible form
    % with the positive-weight convention. The non-repeated basis of periodic
    % tensors is introduced later, at lines 287--303, and is not part of this
    % definition.
    It extends the normal canonical form
    (Definition~\ref{def:is_normal_canonical_form}),
    which requires the stronger condition that every block is $1$-periodic.
\end{definition}

\begin{remark}[Normal canonical form gives irreducible form]
    \label{rem:irr_form_vs_ncf}
\lean{MPSTensor.toIsIrreducibleFormOfPhaseNormalized,
    MPSTensor.toIsIrreducibleFormOfPhaseNormalized_period_eq_one}
\leanok
    Let $A^i=\bigoplus_j \mu_j A_j^i$ be a normal canonical block family with
    nonzero complex weights. Replacing $\mu_j$ by $|\mu_j|$ and $A_j$ by
    $(\mu_j/|\mu_j|)A_j$ gives the positive-weight irreducible form convention
    and does not change the MPV family. The corresponding irreducible-form
    witness has all periods equal to~$1$. This is the period-one part of the statement in
    \cite[Section~II]{DeLasCuevas2017Irreducible} that normal tensors are
    periodic tensors with period~$1$.
\end{remark}

\begin{remark}[Irreducible form allows higher periods]
    \label{rem:irreducible_form_higher_periods}
\notready
    The converse containment is not part of the implication above.
    Irreducible form permits periodic blocks with period larger than~$1$; for
    instance, the antiferromagnetic example has period~$2$.
\end{remark}

\noindent Physical-index rotation was introduced in
Definition~\ref{def:rotate_physical} (Chapter~\ref{ch:symmetry}).
The following compatibility theorems are proved in the present chapter.

\begin{theorem}[Transfer map preserved by physical rotation]%
    \label{thm:transfer_map_rotate_physical}
    \lean{MPSTensor.transferMap_rotatePhysical}
    \leanok
    \uses{def:rotate_physical, def:transfer_map}
    If $u u^\dagger = I$, then the transfer map is invariant under
    physical-index rotation:
    $\mathcal{E}_{A_u} = \mathcal{E}_A$.
\end{theorem}

\begin{proof}\leanok
    \uses{lem:unitary_kraus_mixing}
    This is the unitary freedom of Kraus representations from
    \cite[Theorem~2.1(4)]{Wolf2012Quantum}.
\end{proof}

\begin{theorem}[Left-canonical preserved by physical rotation]%
    \label{thm:left_canonical_rotate_physical}
    \lean{MPSTensor.isLeftCanonical_rotatePhysical}
    \leanok
    \uses{def:rotate_physical}
    If $u u^\dagger = I$ and $A$ is left-canonical, then
    $A_u$ is left-canonical.
\end{theorem}

\begin{proof}\leanok
    \uses{def:rotate_physical}
    Expand $\sum_i B_i^\dagger B_i$ where $B_i = \sum_j u_{ij} A_j$,
    swap sums, apply orthogonality $u^\dagger u = I$, and collapse to
    $\sum_j A_j^\dagger A_j = I$.
\end{proof}

\begin{theorem}[Irreducibility preserved by physical rotation]%
    \label{thm:irreducible_tensor_rotate_physical}
    \lean{MPSTensor.isIrreducibleTensor_rotatePhysical}
    \leanok
    \uses{def:rotate_physical}
    If $u u^\dagger = I$ and $A$ is irreducible, then
    $A_u$ is irreducible.
\end{theorem}

\begin{proof}\leanok
    \uses{def:rotate_physical}
    Contrapositive: any nontrivial invariant projection~$P$ for the
    rotated tensor yields $(1-P)\, A_k\, P = 0$ for all~$k$ via the
    orthogonality of~$u$, contradicting irreducibility of~$A$.
\end{proof}

\begin{theorem}[Periodicity preserved by physical rotation]%
    \label{thm:periodic_rotate_physical}
    \lean{MPSTensor.isPeriodic_rotatePhysical}
    \leanok
    \uses{def:rotate_physical, thm:transfer_map_rotate_physical,
        thm:left_canonical_rotate_physical,
        thm:irreducible_tensor_rotate_physical}
    If $u u^\dagger = I$ and $A$ is periodic with period~$m$, then
    $A_u$ is periodic with the same period.
\end{theorem}

\begin{proof}\leanok
    \uses{thm:transfer_map_rotate_physical,
        thm:left_canonical_rotate_physical,
        thm:irreducible_tensor_rotate_physical}
    Assembles transfer-map invariance, left-canonical preservation, and
    irreducibility preservation.
\end{proof}

\begin{theorem}[Equality of MPV families preserved by physical rotation]%
    \label{thm:same_mpv_rotate_physical}
    \lean{MPSTensor.sameMPV₂_rotatePhysical}
    \leanok
    \uses{def:rotate_physical, def:same_mpv}
    If $\mathcal{V}(A) = \mathcal{V}(B)$, then
    $\mathcal{V}(A_u) = \mathcal{V}(B_u)$.
\end{theorem}

\begin{proof}\leanok
    \uses{def:rotate_physical, def:same_mpv}
    Induction on the word length with a strengthened hypothesis
    carrying an arbitrary prefix matrix, reducing each step to the
    hypothesis via word concatenation.
\end{proof}

\begin{theorem}[Physical rotation distributes over block-diagonal composition]%
    \label{thm:rotate_physical_blocks}
    \lean{MPSTensor.rotatePhysical_toTensorFromBlocks}
    \leanok
    \uses{def:rotate_physical}
    The rotated block-diagonal tensor satisfies
    $(\bigoplus_k \mu_k A_k)_u
        = \bigoplus_k \mu_k (A_k)_u$.
\end{theorem}

\begin{proof}\leanok
    \uses{def:rotate_physical}
    Pointwise matrix-entry computation using linearity of
    the block-diagonal embedding and reindexing.
\end{proof}

\begin{theorem}[Irreducible form preserved by physical rotation]%
    \label{thm:irreducible_form_rotate_physical}
    \lean{MPSTensor.isIrreducibleForm_rotatePhysical}
    \leanok
    \uses{def:rotate_physical, def:irreducible_form,
        thm:periodic_rotate_physical,
        thm:same_mpv_rotate_physical,
        thm:rotate_physical_blocks}
    If $u u^\dagger = I$ and $A$ is in irreducible form, then
    $A_u$ is in irreducible form with the same block structure and
    rotated blocks.
\end{theorem}

\begin{proof}\leanok
    \uses{thm:periodic_rotate_physical,
        thm:same_mpv_rotate_physical,
        thm:rotate_physical_blocks}
    Each block remains periodic by
    Theorem~\ref{thm:periodic_rotate_physical}, equality of MPV families
    transfers via Theorem~\ref{thm:same_mpv_rotate_physical}, and
    block-diagonal compatibility follows from
    Theorem~\ref{thm:rotate_physical_blocks}.
\end{proof}
